\scanpage{16}
Polygône régulier (ou réalisation régulière d'un polygône combinatoire) : définitions possibles (cotés égaux, angles égaux, \inserted{\deleted{\unsure{il avait|il y avait|il suffit}} les deux,} transitivité groupe des symétries ...) Existence, unicité modulo homothétie : tout isom. combinatoire est induit par une unique similitude.

\figure{\begin{tikzpicture}[x=.8cm,y=.8cm]
\draw (0,0) -- (1.14,.66);
\draw (1.18,.70) -- (1.55,1.97);
\end{tikzpicture}}

Notion d'enveloppe affine (ou vectorielle) d'un polygône combinatoire ... \unsure{sens.|jeu.} précisé

Fonctorialité de $\overline{\underline{X\times_G Y\quad\text{p.r. à }(G,X,Y)}}$ ...

\[
\begin{aligned}
\unsure{P|\overline{P}}\times_G E={}^{P}\!\deleted{B}E=F
&\qquad\qquad &P&\quad G\text{-torseur}\\
&&E&\quad G\text{-ensemble}
\end{aligned}
\]

$P\longrightarrow\operatorname{Bij}(E,F)$ compatible aux opérations de $G$ à droite

\uline{NB} Si $G$ opère fidèlement sur $E$, $G$ opère ``sans pts fixes'' sur $\operatorname{Bij}(E,F)$ --- i.e. $\deleted{\unsure{y}}u.g=u$ $(g\in G,\ x\in B)$ implique $g=1$, ou encore les $G_u$ $(u\in B)$ sont réduits à $\{1\}$, ou encore $B$ est une somme de torseurs ... Dans ce cas\inserted{ et dans ce cas seulement}, $P\mathbin{\text{---}}\operatorname{Bij}(E,F)$ est injectif

\[
P\hookrightarrow\operatorname{Bij}(E,F)
\]

% slanted
\margin{$(G,E)$-structure\\sur un $F$\\\deleted{\unsure{Algèbre|Alg.}}}
\noindent$\left[\begin{minipage}{.84\linewidth}
En tous les cas, $F\inserted{={}^{P}\!E}$ est muni de la structure suivante : donnée d'une orbite de $G$ dans $\operatorname{Bij}(E,F)$.
\end{minipage}\right.$

{\footnotesize [On verra que (si $G$ opère $\text{fid}^{\mathrm{t}}$ sur $E$) la donnée \deleted{\unsure{d'un|ens}} d'un \deleted{\unsure{$(G,E)$-st}} $F$ avec $(G,E)$-structure \deleted{\unsure{équiv.|équivaut}} celle \unsure{de|d'un} $G$-torseur \unsure{à équiv. près.|et applic$^{\mathrm{tn}}$}.]}

\uline{Corollaire} \unsure{Point à faire|Point à finir}
% archive sheet number: 15
